% Shared preamble for the Midterm 1 materials (help sheet, study guide,
% practice exams, solutions). \input after \documentclass. Every document
% that uses it starts with \DocumentMetadata and compiles under LuaLaTeX.
\usepackage[letterpaper, margin=0.75in, top=0.55in, bottom=0.6in]{geometry}
\usepackage{fontspec}
\setmainfont{Fira Sans}
\usepackage{amsmath}
% Math in the same family as the text; without this the formulas are set
% in Computer Modern and look pasted in from another document.
\usepackage{unicode-math}
\setmathfont{Fira Math}
\usepackage{graphicx}
\usepackage{booktabs}
\usepackage{array}
\usepackage{multicol}
\usepackage{needspace}
\usepackage{xcolor}
\usepackage[most]{tcolorbox}
\definecolor{accent}{HTML}{BF5700}
\definecolor{accentdark}{HTML}{8F4100}
\definecolor{cream}{HTML}{FDF3EA}
\definecolor{muted}{HTML}{595A5B}
\definecolor{rule}{HTML}{D9D9D9}
\usepackage{titlesec}
\titleformat{\section}{\Large\bfseries\color{accent}}{}{0pt}{}
\titlespacing*{\section}{0pt}{0pt}{4pt}
\titleformat{\subsection}{\normalsize\bfseries\color{accentdark}}{}{0pt}{}
\titlespacing*{\subsection}{0pt}{9pt}{2pt}
% Running header for the help sheet, as on the 340 and 441 sheets.
\usepackage{fancyhdr}
\fancypagestyle{sheet}{%
  \fancyhf{}\renewcommand{\headrulewidth}{0.4pt}%
  \fancyhead[L]{\small ECON 201: Principles of Microeconomics}%
  \fancyhead[R]{\small Midterm 1 Help Sheet}%
  \fancyfoot[C]{\small\thepage}}
\usepackage{hyperref}
\hypersetup{hidelinks}
\setlength{\parindent}{0pt}
\setlength{\parskip}{3pt}
\setlength{\abovedisplayskip}{4pt}
\setlength{\belowdisplayskip}{4pt}
\setlength{\abovedisplayshortskip}{2pt}
\setlength{\belowdisplayshortskip}{4pt}
\setlength{\columnsep}{22pt}

% Compact lists.
\makeatletter
\def\@listi{\leftmargin\leftmargini \topsep 2pt \partopsep 0pt \parsep 0pt \itemsep 1pt}
\let\@listI\@listi
\def\@listii{\leftmargin\leftmarginii \labelwidth\leftmarginii
  \advance\labelwidth-\labelsep \topsep 1pt \parsep 0pt \itemsep 0pt}
\makeatother

% Document header: title at left, course at right, with a rule beneath.
\newcommand{\docheader}[2]{%
  {\LARGE\bfseries\color{accent}#1}\hfill{\large #2}\par
  \vspace{2pt}{\color{accent}\rule{\linewidth}{0.8pt}}\par\vspace{6pt}}

% Panels. Cream for examples and reference material, white with an orange
% frame for things the student fills in. Text on cream uses accentdark
% (7.7:1), since the brand orange falls below 4.5:1 there.
\newtcolorbox{panel}[1]{colback=cream, colframe=accent, boxrule=0.6pt, arc=2pt,
  left=7pt, right=7pt, top=4pt, bottom=5pt, title=#1, fonttitle=\bfseries,
  coltitle=white, colbacktitle=accent, toptitle=2pt, bottomtitle=2pt,
  before skip=6pt, after skip=6pt}
\newtcolorbox{workpanel}[1]{colback=white, colframe=accent, boxrule=0.6pt, arc=2pt,
  left=7pt, right=7pt, top=4pt, bottom=5pt, title=#1, fonttitle=\bfseries,
  coltitle=white, colbacktitle=accent, toptitle=2pt, bottomtitle=2pt,
  before skip=6pt, after skip=6pt}
\newtcolorbox{notepanel}[1]{colback=white, colframe=muted, boxrule=0.5pt, arc=2pt,
  left=7pt, right=7pt, top=3pt, bottom=4pt, title=#1, fonttitle=\bfseries,
  coltitle=white, colbacktitle=muted, toptitle=1pt, bottomtitle=1pt,
  before skip=6pt, after skip=6pt}

% Pointers back to the course materials.
\newcommand{\where}[1]{{\color{muted}\small #1}}
% An empty box to tick by hand.
\newcommand{\tick}{{\large$\mdlgwhtsquare$}}
% Step label inside a worked example.
% Help sheet rows: the formula, then what its symbols mean. The argument
% is the width of the formula column; both cells are vertically centred so
% a note sits level with a tall fraction.
\newenvironment{formulas}[1]
  {\begin{tabular}{@{}>{\raggedright\arraybackslash$\displaystyle}m{#1}<{$}@{\hspace{10pt}}>{\small\color{muted}\raggedright\arraybackslash}m{\dimexpr\linewidth-#1-10pt\relax}@{}}}
  {\end{tabular}\par\vspace{3pt}}
\newcommand{\frow}{\\[5pt]}
% One formula on its own line, left-aligned, for the help sheet.
\newlength{\flskip}\setlength{\flskip}{5pt}
\newcommand{\fl}[1]{\par\noindent\hspace*{1em}$\displaystyle #1$\par\vspace{\flskip}}
% Exam pieces. \pts marks a part's points; \work leaves ruled space on an
% exam and nothing in its solutions; \sol prints only in the solutions.
\newif\ifsolutions
\def\ptsone{1}
\DeclareRobustCommand{\pts}[1]{{\def\ptsarg{#1}\color{muted}\small[#1~\ifx\ptsarg\ptsone pt\else pts\fi]}}
\pdfstringdefDisableCommands{\def\pts#1{}}
% A part of a problem: label and points, with air above it.
% A hint or aside under a part, in the muted colour.
\newcommand{\fnote}[1]{\par\vspace{2pt}{\small\color{muted}#1}\par}
% A physical sheet of the exam: starts on an odd page (a blank page is
% inserted if needed, so sheets can be separated and graded by different
% people) and carries a name and CWID line on its front.
% The CWID is nine boxes, one digit each, as on the scanned exam.
\newcommand{\cwidbox}{\fbox{\rule{0pt}{7mm}\hspace{5.5mm}}\hspace{0.6mm}}
\newcommand{\cwidboxes}{{\setlength{\fboxsep}{0pt}\setlength{\fboxrule}{0.5pt}%
  \cwidbox\cwidbox\cwidbox\cwidbox\cwidbox\cwidbox\cwidbox\cwidbox\cwidbox}}
\newcommand{\sheet}[1]{\clearpage\ifodd\value{page}\else\mbox{}\thispagestyle{empty}\clearpage\fi
  \noindent Name: \rule{7.5cm}{0.4pt}\par\vspace{8pt}
  \noindent CWID:\enspace\raisebox{-2.4mm}{\cwidboxes}\enspace{\small\color{muted}One digit per box.}\par\vspace{5pt}
  \noindent{\small\color{muted}#1}\par\vspace{6pt}}
% Multiple choice laid out like the scanned exam: the question and its
% options at left, a row of answer bubbles at right (exam copy only). The
% bubble is Fira Math's large circle; Fira Sans has none.
\newfontfamily\bubblefont{Fira Math}
\newcommand{\bubble}[1]{\shortstack{{\bubblefont\large ◯}\\[-1pt]{\scriptsize\color{muted}#1}}}
\newcommand{\mcopt}[2]{\par\noindent\hangindent=1.8em\hangafter=1\makebox[1.8em][l]{(#1)}#2}
\newcommand{\mcq}[5]{\item\ifsolutions\begin{minipage}[t]{\linewidth}\else\begin{minipage}[t]{0.74\linewidth}\fi
  #1\par\vspace{3pt}\mcopt{A}{#2}\mcopt{B}{#3}\mcopt{C}{#4}\mcopt{D}{#5}\end{minipage}%
  \ifsolutions\else\hfill\begin{minipage}[t]{0.22\linewidth}\raggedleft
  \bubble{A}\hspace{5pt}\bubble{B}\hspace{5pt}\bubble{C}\hspace{5pt}\bubble{D}\end{minipage}\fi\par}
\newcommand{\qpart}[2]{\par\vspace{10pt}\noindent(#1) \pts{#2} }
% A numbered short-answer question: number and points, with more air above
% it than a part.
\newcommand{\sa}[2]{\par\vspace{14pt}\noindent\textbf{#1.} \pts{#2} }
\newcount\answerlinecount
\newcommand{\answerlines}[1]{%
  \par\vspace{1pt}\answerlinecount=#1
  \loop\ifnum\answerlinecount>0
    \rule{\linewidth}{0.3pt}\par\vspace{6pt}%
    \advance\answerlinecount by -1
  \repeat
  \vspace{-6pt}}
% Blank working space on an exam, about #1 lines of handwriting; nothing in
% the solutions. Ruled lines get in the way of fractions and arithmetic.
\newlength{\workline}\setlength{\workline}{18pt}
\newcommand{\work}[1]{\ifsolutions\else\par\vspace{#1\workline}\fi}
\newcommand{\sol}[1]{\ifsolutions\par\vspace{1pt}{\color{accentdark}\textbf{Solution.} #1}\par\fi}
\newcommand{\blankrow}{\rule{0pt}{1.9em}}
% Percentage change, as in the slides.
\newcommand{\pc}{\%\,\Delta}
\newcommand{\step}[1]{\textbf{\color{accentdark}#1}\enspace}
